% STATUS: DRAFT
\chapter{Normal states, weights, and traces}\label{ch:traces}

\begin{physmotiv}
The entropy $-\Tr\rho\log\rho$ requires two things: a density matrix $\rho$ and a trace $\Tr$ to take its expectation. In $\BH$ both exist. In a general von Neumann algebra we must ask which of them survive. The answer is a hierarchy: \emph{normal states} always exist and are the physical states; a \emph{trace} may or may not exist; and when it does not, one still has \emph{weights}, which are traces' unbounded, non-tracial cousins, and which will be the main characters in modular theory and in crossed products.
\end{physmotiv}

\section{Normal states}

A state $\omega$ on a von Neumann algebra $\M$ is \emph{normal} if it is $\sigma$-weakly continuous; equivalently (\cref{ch:doublecomm}) it is completely additive on orthogonal projections. A state is \emph{faithful} if $\omega(a\ad a)=0\Rightarrow a=0$, and its \emph{support} is the smallest projection $s$ with $\omega(s)=1$. In a representation, a vector $\xi$ defines the normal \emph{vector state} $\omega_\xi(a)=\inner\xi{a\xi}$. For $\M=\BH$:
\begin{center}
normal state $\Longleftrightarrow$ $\omega=\Tr(\rho\,\cdot)$ with $\rho\ge0$ trace class, $\Tr\rho=1$.
\end{center}
For general $\M$ normal states play the role of density matrices without the density: there may be no operator $\rho$ in $\M$ at all. Key results:
\begin{itemize}
\item $\omega$ is faithful and normal on $\M$ iff $\omega=\omega_\xi$ for a vector $\xi$ that is \emph{separating} for $\M$ (then $\xi$ is cyclic for $\M'$), and every normal state is a vector state in a suitable representation (the \emph{standard form}, \cref{ch:tomita}).
\item Every normal state on $\M$ is a vector state in the amplification $\M\otimes\id_{\ell^2}$ (purification): a statement of the existence of purifications.
\end{itemize}

\section{Traces}

\begin{definition}
A \emph{trace} on $\M$ is a map $\tau:\M_+\to[0,\infty]$ that is additive, positively homogeneous and satisfies $\tau(x\ad x)=\tau(xx\ad)$ for all $x\in\M$ (so $\tau(uau\ad)=\tau(a)$ for unitaries). It is \emph{faithful} if $\tau(a)=0\Rightarrow a=0$, \emph{semifinite} if every non-zero projection dominates a non-zero projection of finite trace, and \emph{normal} if $\tau(\sup a_i)=\sup\tau(a_i)$ for bounded increasing nets. We write ``f.n.s.\ trace'' for faithful normal semifinite.
\end{definition}

On $\BH$ the usual $\Tr$ is the unique f.n.s.\ trace up to scale. Several things are worth noting:
\begin{itemize}
\item If $\tau(\id)<\infty$, $\tau$ extends to a linear tracial functional, which we can normalize into a \emph{tracial state}; for $M_n$ it is $\Tr/n$, for the spin chain $\tau=\bigotimes\frac12\Tr$.
\item For a \emph{factor}, f.n.s.\ traces are unique up to scale if they exist. For non-factors there is one scale per sector.
\item \emph{Not every von Neumann algebra has a (semifinite) trace}: the algebras of QFT do not (type III, \cref{ch:classification}).
\end{itemize}

\subsection{Density matrices relative to a trace}
Suppose $\M$ has a f.n.s.\ trace $\tau$. Then for every normal state $\omega$ there is a unique positive operator $\rho_\omega$ \emph{affiliated} with $\M$ (\cref{ch:spectral}) with $\tau(\rho_\omega)=1$ and
\begin{equation}
\omega(a)=\tau(\rho_\omega a)\qquad(a\in\M).
\end{equation}
This is a Radon--Nikodym theorem; $\tau$ plays the role of the reference measure. The \emph{von Neumann entropy} is then
\begin{equation}
S(\omega)=-\tau(\rho_\omega\log\rho_\omega).
\end{equation}
The entropy depends on the normalization of $\tau$: rescaling $\tau\to c\tau$ shifts $S$ by $\log c$ (a state-independent constant), exactly as one expects. If $\M$ is a factor of type II$_1$, with the trace normalized to $\tau(\id)=1$, then $S(\omega)\le0$ with the maximum $0$ attained for $\rho=\id$, the tracial state. In general:
\begin{center}
\emph{entropy is defined on $\M$ iff $\M$ has a f.n.s.\ trace, i.e.\ iff $\M$ is of type I or II.}
\end{center}

\section{Weights}

When no trace exists, one needs a weaker structure.

\begin{definition}
A \emph{weight} on $\M$ is an additive, positively homogeneous map $\varphi:\M_+\to[0,\infty]$. Weights are normal, faithful, semifinite under the same conditions as above; a trace is a weight with $\varphi(x\ad x)=\varphi(xx\ad)$.
\end{definition}
Normal states are the bounded (finite) weights ($\varphi(\id)=1$). The theorem of Haagerup states that every normal weight is a supremum of normal positive functionals, and every von Neumann algebra (on a separable space) has a faithful normal semifinite weight. For example:
\begin{itemize}
\item On $\BH$: $\varphi(a)=\Tr(\rho a)$ for any positive unbounded $\rho$ (e.g.\ $\rho=\ee^{-\beta H}$ with $\Tr\ee^{-\beta H}=\infty$, as in a system with continuous spectrum). Faithful normal semifinite iff $\rho$ is non-singular.
\item On a type III algebra: faithful normal states $\omega_\xi$ for $\xi$ cyclic and separating (the vacuum for a wedge). These are weights, not traces, and they are \emph{not} ``$\Tr\rho\,\cdot$'' for any $\rho$ in $\M$ because there is no trace.
\end{itemize}
Crucially, each f.n.s.\ weight $\varphi$ defines a modular automorphism group $\sigma_t^\varphi$ (\cref{ch:tomita}), and the weight is a trace iff $\sigma^\varphi$ is trivial. This is the sense in which modular theory measures the failure of the trace.

\begin{whatwrong}
Do not confuse the following: ``$\omega$ is faithful'' (a statement about a state), ``$\M$ has a faithful normal state'' (true for all $\sigma$-finite algebras), and ``$\M$ has a faithful normal \emph{trace}'' (a property of the algebra, true only for types I and II, in other words a restriction on its type). The first two hold in QFT; the third does not.
\end{whatwrong}

\begin{keyidea}
Normal state: physical state, always available. Trace: needed for density matrices and entropy, available only in types I, II. Weights: unbounded traces' generalization, always available and the entry point to modular theory.
\end{keyidea}

\section*{Exercises}
\begin{exercise}
Show that if $\tau$ is a tracial state on a factor $\M$ and $\omega$ is a faithful normal state, then $\rho_\omega$ has trivial kernel (it need not be boundedly invertible in general, but it is in $M_n$). Compute $S(\omega)$ for $\M=M_n$, and for $\M=M_n\otimes M_k$ with a product state.
\end{exercise}
\begin{exercise}
Let $\M$ be II$_1$ with $\tau(\id)=1$. Using Jensen's inequality for the convex function $x\log x$ prove $S(\omega)\le0=S(\tau)$ for every normal state $\omega=\tau(\rho\,\cdot)$.
\end{exercise}
\begin{exercise}
Let $\varphi=\Tr(\ee^{-\beta H}\,\cdot)$ on $\BH$ with $H$ having spectrum $[0,\infty)$ with Lebesgue density of states. Show $\varphi$ is a faithful normal semifinite weight, is not a state, and that $\sigma^\varphi_t(a)=\ee^{-\ii\beta Ht}a\,\ee^{\ii\beta Ht}$ is not trivial (hence $\varphi$ is not a trace). (We will confirm this in \cref{ch:tomita}.)
\end{exercise}
\begin{exercise}
Show that if $\tau_1,\tau_2$ are two f.n.s.\ traces on a factor then $\tau_2=c\tau_1$ (hint: consider the Radon--Nikodym derivative $h$, which must be central).
\end{exercise}
